% Generated draft: requires mathematical and source review.
\documentclass[11pt]{article}
\usepackage[margin=25mm]{geometry}
\usepackage{fontspec}
\setmainfont{Latin Modern Roman}
\usepackage{amsmath,amssymb,mathtools,unicode-math}
\setmathfont{Latin Modern Math}
\usepackage{enumitem,needspace,xurl,hyperref}
\setlength{\parindent}{0pt}
\setlength{\parskip}{5pt}
\setlength{\emergencystretch}{3em}
\newcommand{\N}{\mathbb{N}}
\newcommand{\Z}{\mathbb{Z}}
\newcommand{\Q}{\mathbb{Q}}
\newcommand{\R}{\mathbb{R}}
\newcommand{\C}{\mathbb{C}}
\newcommand{\F}{\mathbb{F}}
\newcommand{\A}{\mathbb{A}}
\newcommand{\PP}{\mathbb{P}}
\newcommand{\E}{\mathbb{E}}
\newcommand{\eps}{\varepsilon}
\newcommand{\dd}{\,\mathrm{d}}
\begin{document}
\section*{UnsolvedMath ID: 20002017}
\textbf{Source problem code:} catalogue entry
\subsection{Definitions and mathematical statement}
Thomas Branson's Problem 1a asks what information about a metric functional, already understood along conformal changes, can be obtained by moving across conformal classes. For an integral of a local invariant it suggests computing an anti-conformal variation; for a nonlocal spectral invariant such as a functional determinant, computing that variation is itself part of the challenge. The source does not fix a functional, operator, dimension, or precise meaning of anti-conformal.

The supplied aid adopts the following infinitesimal interpretation. Let \(M\) be a closed smooth manifold of dimension \(n\), \(g\) a Riemannian metric, and \(F\) a differentiable functional having a smooth \(L^2\) gradient \(T_g\), so that for every smooth symmetric covariant two-tensor \(h\),
\[
 D F_g(h)=\int_M\langle T_g,h\rangle_g\dd V_g.
\]
Infinitesimal conformal changes are \(h=2ug\), \(u\in C^\infty(M)\), and their \(L^2\)-orthogonal complement consists of trace-free tensors. Write
\[
 T_g^0=T_g-\tfrac1n(\operatorname{tr}_gT_g)g.
\]
Then the reported two responses are
\[
 D F_g(2ug)=2\int_Mu\operatorname{tr}_gT_g\dd V_g,
 \qquad D F_g(h)=\int_M\langle T_g^0,h\rangle\dd V_g
 \quad(\operatorname{tr}_gh=0).
\]
Thus conformal information records the trace of the gradient, and the chosen anti-conformal information records its trace-free part. This is one working interpretation of the workshop wording.

If \(F\) is invariant under diffeomorphisms, the aid reports \(\operatorname{div}_gT_g=0\). A transverse-traceless tensor is a smooth \(h\) with both \(\operatorname{tr}_gh=0\) and \(\operatorname{div}_gh=0\). Its no-loss test states
\[
 \bigl[D F_g(2ug)=0\ \forall u,
 \quad D F_g(h)=0\ \forall h\text{ transverse-traceless}\bigr]
 \quad\Longleftrightarrow\quad D F_g(h)=0\ \forall h.
\]
The hypotheses of diffeomorphism invariance and a smooth \(L^2\) gradient remain essential. The assertion does not identify all trace-free tensors with transverse-traceless tensors.

For the Einstein--Hilbert functional \(F(g)=\int_M R_g\dd V_g\) in dimension \(n\ge3\), using covariant metric variation, the aid specializes to
\[
 T_g=-\operatorname{Ric}_g+\tfrac12R_gg,\qquad
 T_g^0=-\operatorname{Ric}_g^0.
\]
Consequently vanishing on every trace-free variation is equivalent to \(\operatorname{Ric}_g^0=0\), and the trace-free direction \(h=\operatorname{Ric}_g^0\) gives
\[
 D F_g(\operatorname{Ric}_g^0)
 =-\|\operatorname{Ric}_g^0\|_{L^2(g)}^2<0
 \quad\text{if }\operatorname{Ric}_g^0\not\equiv0.
\]
This is a trace-free descent certificate. It is not automatically a transverse-traceless direction: in these conventions
\(\operatorname{div}_g\operatorname{Ric}_g^0=(\tfrac12-\tfrac1n)\dd R_g\).
The aid supplies no anti-conformal determinant formula; nonlocal spectral regularity, operator choice, zero modes, and determinant conventions remain to be specified before that part of the research request can be made precise.
\subsection{Short English statement}
The source asks how metric functionals change across conformal classes. The aid interprets these directions as trace-free variations and separates the gradient's trace and trace-free responses. For diffeomorphism-invariant functionals it gives a full-stationarity test using conformal and transverse-traceless directions, while its Einstein--Hilbert descent direction is trace-free and need not itself be transverse-traceless.
\subsection{Sources}
\begin{itemize}
\item Local UnsolvedMath registry, permanent ID 20002017, exact source code catalogue entry; canonical \texttt{aim-geometry-notes.json}, \texttt{notes[354]}, tag \texttt{problem}; AIM Workshop Problem Lists, Geometry. Created and updated 14 August 2026.
\item AIM workshop \emph{Conformal structure in geometry, analysis, and physics}, 12--16 August 2003, subsection ``Thomas Branson. Anti-conformal perturbations'', Problem 1a; \url{https://aimath.org/WWN/confstruct/confstruct.pdf}. This attribution and location are supplied by the verification; the PDF was not accessed.
\item Attempt 1: \texttt{partial\_result}, problem status at run \texttt{partially\_solved}, rigor label \texttt{proved}, novelty confidence \texttt{low}; classification \texttt{PARTIAL-PROGRESS}, registry status \texttt{partially\_solved}. Candidate: conformal/transverse-traceless stationarity lemma and Einstein--Hilbert trace-free descent certificate. No separate authoritative theorem citation is supplied.
\item Ulam AI / UnsolvedMath Contributors; CC BY 4.0. Machine-generated research aid, not peer reviewed; independently verify before citation.
\end{itemize}
\end{document}
